\section{Backbone, Information Flow, and Design Confirmation}
\label{app:q3}

\subsection{Parameter Organization}
The backbone comparison holds history $H$, future representation $Z$, and
reference information flow $G_0$ fixed. In $G_0$, future supervision updates
shared parameters, but Action does not read a generated future. D shares
attention QKV/output projections across observation/meta, future, and Action,
with separate feed-forward networks. Adapted E has two parameter groups,
World and Action, each with its own QKV, output projection, adaptive layer
normalization (AdaLN), and feed-forward network. Observation/meta and active
future tokens share World parameters. Observation/meta tokens use diffusion
time $\tau=0$; generated future tokens retain their own time. When the future
is only a deterministic readout, neither backbone instantiates an unused
generative stream.

Internal normalization is part of this comparison. E's World AdaLN has
randomly initialized scale/shift biases at $\tau=0$, whereas D's observation
LayerNorm starts with affine scale one and shift zero. D/E therefore compares
parameter organization and normalization as complete backbone choices.
The early E system, with different inputs and a joint program, remains a
separate system reference.

Both arms use fresh training seeds 0, 1, and 2. Common input projections,
history/time parameters, and readouts start identically within a seed;
backbone groups use separate initialization streams. The reference recipe
is four epochs, global batch 16, 40,104 updates, and 641,664 consumed chunks
per run (Appendix~\ref{app:implementation}). This gives six endpoints and
8,400 development rollouts. D uses FFN width $4d$, with $d=768$.
E chooses its FFN width from positive multiples of 64 to minimize the
absolute active-backbone parameter difference, breaking ties toward the
smaller width. The difference must be within 5\%; depth, attention width,
and heads remain fixed. Frontend, teacher, and total parameters are listed separately.
Equal exposure does not imply equal FLOPs.

\paragraph{Selection rule.}
The fixed reference device is the previously measured NVML-identified
\texttt{NVIDIA RPBZZZ6}, capacity 102,641,958,912 bytes, with FP32/TF32.
All candidates use this same device and precision, with p95 time-to-action
at most 200\,ms and peak memory at most 80\% of device capacity. For candidate
minus reference SR, a 95\% paired interval wholly above $+2$ percentage points
supports a practically meaningful advantage; an interval inside $[-2,2]$
supports the specified equivalence criterion. A candidate with at least 10\%
lower p95 latency can be selected on cost if its SR interval lies above
$-2$ points. Per-seed scores and native task/safety outcomes accompany the
decision. If three seeds do not resolve it, both arms add seeds 3 and 4.
If neither dominates after this extension, the reference is retained without
a superiority claim. The default complete reference is T8 RGB with causal
estimated motion, scene tracks plus target-pose supervision, D, and R1
action-only deployment. Oracle-information studies retain their separate
input track. Table~\ref{tab:backbone_confirmation} records the backbone
comparison and its revisit under the selected information flow.

% Planned results: em-dash cells are pending measurements, not zeros.
\begin{table*}[t]
\centering\small
\caption{\tablelang{Backbone comparison and directed revisit. SR averages common development scenes and training seeds. Parameter counts cover the active backbone; p95 is time to the first issued action. The revisit fixes the selected information flow.}{主干比较及定向复验。SR汇总共同开发场景及训练种子；参数量仅计活动主干，p95为首条动作前的策略耗时。复验固定在开发集选出的信息流。}}
\label{tab:backbone_confirmation}
\begin{tabular}{@{}llrrrrr@{}}
\toprule
\tablelang{Flow}{信息流}&\tablelang{Backbone}{主干}&\tablelang{Params (M)}{参数量(M)}&SR (\%)&$\Delta$SR (pp)&p95 (ms)&\tablelang{Memory (GiB)}{显存(GiB)}\\
\midrule
$G_0$&D&\resultplaceholder&\resultplaceholder&\tablelang{ref.}{参照}&\resultplaceholder&\resultplaceholder\\
$G_0$&E&\resultplaceholder&\resultplaceholder&\resultplaceholder&\resultplaceholder&\resultplaceholder\\
$G^\star$&D&\resultplaceholder&\resultplaceholder&\tablelang{ref.}{参照}&\resultplaceholder&\resultplaceholder\\
$G^\star$&E&\resultplaceholder&\resultplaceholder&\resultplaceholder&\resultplaceholder&\resultplaceholder\\
\bottomrule
\end{tabular}
\end{table*}


\subsection{A Common Generated-Future Interface}
\label{app:generated_future}
Every information-flow row binds the same generated state: scene
displacement $F\in\mathbb R^{196\times8\times3}$, or sparse image tracks
$F\in\mathbb R^{128\times8\times2}$, over 50--400\,ms. Displacements use the
fixed robot-base frame. Tracks retain current source identities and use
image coordinates divided by width and height. Source coordinates and
current validity are observation inputs. Future values are standardized
per coordinate channel over valid training elements, with a minimum scale
of $10^{-6}$; all rows use the same statistics. Each source point's full
trajectory is projected to one token, with source-location and representation
embeddings; the output projection reconstructs its eight ordered horizons.

For standardized trajectory $x_1$ and $\epsilon\sim\mathcal N(0,I)$,
$x_\tau=(1-\tau)\epsilon+\tau x_1$ has FM target $x_1-\epsilon$.
One $\tau\sim U[0,1)$ is shared within each future tensor. Coordinate MSE
averages valid elements. Invalid future coordinates retain sampled source
noise throughout interpolation and are excluded from loss;
their masks do not condition the model. Current source padding is masked
identically at training and deployment. A separate logit head predicts
track visibility with BCE on valid correspondences. Ground-truth future
visibility never enters the generator or Action. Target pose remains a
deterministic meta-token readout with the geometric loss in
Appendix~\ref{app:implementation}; it is not the generated state $F$.
The World-to-Action edge exposes future tokens and their hidden states.

This normalized Gaussian interface is common to the information-flow
comparison. The earlier representation screen used a 0.02\,m metric
scene-flow source and deterministic track readouts, so its checkpoints
do not replace matched information-flow training. Auxiliary weights are
fixed across R1--R6 and Alpha is zero.

\subsection{Seven Information-Flow Programs}
\begin{table*}[t]
\centering\small
\caption{Planned Q3 conditions; no measured result is assigned to these rows.
F denotes a generated future state and A an action state. Arrows specify
information access, independently of whether the future loss is enabled.
R1 has two deployments of the same trained checkpoint. Sequential second
stages use clean ground truth in training and generated first-stage values at test time.}
\label{tab:q3_plan}
\begin{tabularx}{\textwidth}{@{}lcccXX@{}}
\toprule
ID & Future loss & $F\to A$ & $A\to F$ & Cross-stream training input & Inference order \\
\midrule
R0 & No & --- & --- & None & A only \\
R1 & Yes & --- & --- & None & A only \\
R1 & Yes & --- & --- & None & Joint F and A \\
R2 & Yes & Yes & --- & Synchronous noisy states & Joint F and A \\
R3 & Yes & --- & Yes & Synchronous noisy states & Joint F and A \\
R4 & Yes & Yes & Yes & Synchronous noisy states & Joint F and A \\
R5 & Yes & Yes & --- & Clean F for second-stage A & F, then A \\
R6 & Yes & --- & Yes & Clean A for second-stage F & A, then F \\
\bottomrule
\end{tabularx}
\end{table*}

Let $C$ denote observation/meta context, $F$ generated future tokens, and
$A$ the $8\times7$ action state. Context is causal in capture time and cannot
read $F$ or $A$. Each generated stream reads $C$ and itself;
Table~\ref{tab:q3_plan} gives the additional cross-stream edges. Future
labels therefore cannot reach Action indirectly through context. A blocked
edge has attention logit $-\infty$ at every layer and head in training and
testing.

\paragraph{Synchronous programs.}
R0 removes future loss and tokens, preserving the action/context interface.
R1 adds future supervision with both cross-stream edges blocked; its same
checkpoint runs action-only or generates an unused future. R2 enables
$F\!\to\!A$, R3 enables $A\!\to\!F$, and R4 enables both. Training uses
independent Gaussian sources and independent uniform diffusion times for
$F$ and $A$, with per-stream time embeddings. At deployment both times
advance on $\{0,.2,.4,.6,.8\}$: five combined evaluations. Streams exchange
their noisy states and hidden representations, not clean future labels.

\paragraph{Sequential programs.}
R5 first generates $F$ with Action absent, then holds its decoded tokens fixed
while generating $A$ through $F\!\to\!A$. Training uses two passes: future FM
with Action absent, then action FM conditioned on clean ground-truth future
tokens. The second pass does not backpropagate through a sampled first
pass; both losses update shared parameters normally. R6 reverses these
roles: action FM, then future FM conditioned on clean actions. At deployment
the second stage reads first-stage predictions. Each stage uses five Euler
steps, totaling ten evaluations. The training/deployment conditioning
difference belongs to the sequential program. R6 can issue its first-stage
action immediately; both time-to-action and time-to-complete-forecast are
reported.

Three seeds give 21 trained endpoints and 24 deployment conditions, including
R1's second mode: 33,600 development rollouts. Every row issues one command
per replan with the same physical limits. World/Action NFE, forward calls,
active parameters, exposure, p50/p95 time-to-action, memory, and physical
forecast errors accompany SR. R0/R1 tests supervision; R1's two deployments
isolate unused generation; R1--R4 tests access; R2/R5 and R3/R6 compare
complete synchronous and sequential programs.

\begin{table*}[t]
\centering\small
\caption{\tablelang{Information-flow results. Each condition uses three training seeds and the same history/future interface. NFE counts are per modality; joint rows share calls. Forecast error is decoded to mm for 3D displacement or pixels for tracks, fixing representation within a comparison. An absent generated output is marked n/a.}{信息流结果。每个条件使用三个训练种子及共同历史/未来接口。NFE按模态计，同步行共享调用。预测误差在物理解码后计算（三维位移为mm，轨迹为像素），同一比较中表示固定。n/a表示没有对应生成输出。}}
\label{tab:information_flow_results}
\begin{tabular}{@{}llccrrrr@{}}
\toprule
ID&\tablelang{Deployment}{部署}&$\mathrm{NFE}_F$&$\mathrm{NFE}_A$&SR (\%)&\tablelang{Forecast error}{预测误差}&p95 (ms)&\tablelang{Safety (\%)}{安全失败(\%)}\\
\midrule
R0&\tablelang{Action}{动作}&0&5&\resultplaceholder&n/a&\resultplaceholder&\resultplaceholder\\
R1&\tablelang{Action}{动作}&0&5&\resultplaceholder&n/a&\resultplaceholder&\resultplaceholder\\
R1&\tablelang{Joint, unread future}{同步，未来不回读}&5&5&\resultplaceholder&\resultplaceholder&\resultplaceholder&\resultplaceholder\\
R2&\tablelang{Joint $F\to A$}{同步$F\to A$}&5&5&\resultplaceholder&\resultplaceholder&\resultplaceholder&\resultplaceholder\\
R3&\tablelang{Joint $A\to F$}{同步$A\to F$}&5&5&\resultplaceholder&\resultplaceholder&\resultplaceholder&\resultplaceholder\\
R4&\tablelang{Joint, bidirectional}{同步双向}&5&5&\resultplaceholder&\resultplaceholder&\resultplaceholder&\resultplaceholder\\
R5&\tablelang{$F$ then $A$}{先$F$后$A$}&5&5&\resultplaceholder&\resultplaceholder&\resultplaceholder&\resultplaceholder\\
R6&\tablelang{$A$ then $F$}{先$A$后$F$}&5&5&\resultplaceholder&\resultplaceholder&\resultplaceholder&\resultplaceholder\\
\bottomrule
\end{tabular}
\end{table*}


\subsection{Interventions and Interaction Tests}
\label{app:design_crosschecks}
Future interventions hold current perception, robot history, sampling arrays,
and decoding fixed. Identity replacement copies native future states.
Zero-motion replacement uses zero 3D displacement or repeats each current
image point across horizons. A matched-current shuffle samples a different
development trajectory with seed 17 in the same task,
object, and visibility stratum, within 5\,mm and $5^\circ$ of current target
pose. Donor sources also match receiver coordinates within two pixels and
have the same role. Both members retain the same valid point/horizon
domain; unmatched windows are excluded with coverage reported. Time shifts use
$\delta_h\pm50$\,ms where labels exist. Noise uses standard deviations
$0,0.5,1,2$ times development prediction RMSE, with shared random draws and
the representation's physical coordinates.

An oracle diagnostic uses measured exogenous target motion before first
contact. Post-contact analysis forks a simulator snapshot under a fixed
candidate action sequence and the same external motion. It measures response
to that candidate's future; an action changed by the intervention need not
produce that same future. R2/R5 model demonstration futures without action
conditioning, whereas R3/R4/R6 can condition on action. Geometric validity
alone does not certify action compatibility. Task and safety outcomes
accompany first-issued translation and SO(3) action changes.

For R2/R4, a clean replacement is standardized using the original training
statistics and inserted at every Euler step as
$F_u^{I}=(1-u)\epsilon_F+uF^{I}$, retaining the native source-noise array.
World hidden states are recomputed before Action is evaluated; R4 keeps
its current Action-to-World edge. Identity controls instead copy native
intermediate $F_u$ states and test exact reconstruction of the unmodified
program. R5 replaces and re-embeds the clean decoded future at its second
stage. R0/R1 and R3/R6, whose issued actions do not read future states,
serve as negative controls. Label validity defines paired analysis support
but is not an input mask. Table~\ref{tab:future_intervention_results}
reports action, contact, task, and coverage measures.

% Planned results: em-dash cells are pending measurements, not zeros.
\begin{table*}[t]
\centering\small
\caption{\tablelang{Frozen future-interface interventions. Action changes compare first-issued commands under matched noise. Contact errors use actual contact events; non-contact failures stay in SR and are excluded only from conditional contact metrics. Coverage is the fraction of eligible windows/scenes with paired intervention support.}{冻结未来接口干预。动作变化比较相同初噪下的首条下发命令。接触误差使用真实接触事件；无接触失败保留在SR中，只不进入条件接触指标。覆盖率为可进行配对干预的有效窗口/场景比例。}}
\label{tab:future_intervention_results}
\begin{tabular}{@{}lrrrrrr@{}}
\toprule
\tablelang{Intervention}{干预}&$\Delta a$ (mm)&$\Delta R$ ($^\circ$)&\tablelang{Contact (ms)}{接触(ms)}&\tablelang{Pose (mm)}{位姿(mm)}&SR (\%)&\tablelang{Coverage (\%)}{覆盖(\%)}\\
\midrule
\tablelang{Native future}{原生未来}&0&0&\resultplaceholder&\resultplaceholder&\resultplaceholder&\resultplaceholder\\
\tablelang{Identity copy}{恒等复制}&\resultplaceholder&\resultplaceholder&\resultplaceholder&\resultplaceholder&\resultplaceholder&\resultplaceholder\\
\tablelang{Zero motion}{零运动}&\resultplaceholder&\resultplaceholder&\resultplaceholder&\resultplaceholder&\resultplaceholder&\resultplaceholder\\
\tablelang{Matched shuffle}{匹配交换}&\resultplaceholder&\resultplaceholder&\resultplaceholder&\resultplaceholder&\resultplaceholder&\resultplaceholder\\
\tablelang{Time offset $\pm50$ ms}{时间偏移$\pm50$ms}&\resultplaceholder&\resultplaceholder&\resultplaceholder&\resultplaceholder&\resultplaceholder&\resultplaceholder\\
\tablelang{Noise $0.5/1/2\times$ RMSE}{噪声$0.5/1/2\times$RMSE}&\resultplaceholder&\resultplaceholder&\resultplaceholder&\resultplaceholder&\resultplaceholder&\resultplaceholder\\
\tablelang{Pre-contact oracle}{接触前oracle}&\resultplaceholder&\resultplaceholder&\resultplaceholder&\resultplaceholder&\resultplaceholder&\resultplaceholder\\
\tablelang{Candidate-conditioned oracle}{候选动作条件oracle}&\resultplaceholder&\resultplaceholder&\resultplaceholder&\resultplaceholder&\resultplaceholder&\resultplaceholder\\
\bottomrule
\end{tabular}
\end{table*}


The confirmation matrix contains three interactions. First, current-only
versus T8 RGB is crossed with current target pose versus pose plus interval
motion, using the selected future target. Second, scene flow versus sparse
tracks is crossed with $G_0$ versus the selected flow, retaining target-pose
supervision. Third, D and E are revisited under that flow. Every cell uses
seeds 0--2 and common development scenes; identical endpoints are reused.
A ranking reversal is reported as an interaction and changes selection of
the complete configuration. The selected $(H,Z,B,G)$ tuple, its action-only
ablation, and the strongest complete alternative are evaluated on the
1,400 independent confirmation scenes in Appendix~\ref{app:protocol},
without further selection.

\begin{table*}[t]
\centering\small
\caption{\tablelang{Targeted interactions and independent confirmation. History rows fix the future design; future/flow rows retain target-pose supervision. Development contrasts match seeds and scenes. ID confirmation uses 1,400 new scenes after configuration selection.}{定向交互与独立确认。历史部分固定未来设计；未来/信息流部分保持目标位姿监督。开发比较匹配种子与场景；ID确认在配置选定后使用新的1,400场景。}}
\label{tab:design_crosschecks}
\begin{tabularx}{\textwidth}{@{}lXrrr@{}}
\toprule
\tablelang{Comparison}{比较}&\tablelang{Condition}{条件}&SR (\%)&$\Delta$SR (pp)&p95 (ms)\\
\midrule
$H_{\rm RGB}\times H_{\rm motion}$&\tablelang{Current RGB / current pose}{当前RGB/当前位姿}&\resultplaceholder&\tablelang{ref.}{参照}&\resultplaceholder\\
&\tablelang{T8 RGB / current pose}{T8 RGB/当前位姿}&\resultplaceholder&\resultplaceholder&\resultplaceholder\\
&\tablelang{Current RGB / interval motion}{当前RGB/区间运动}&\resultplaceholder&\resultplaceholder&\resultplaceholder\\
&\tablelang{T8 RGB / interval motion}{T8 RGB/区间运动}&\resultplaceholder&\resultplaceholder&\resultplaceholder\\
\midrule
$Z\times G$&\tablelang{Scene flow / $G_0$}{场景流/$G_0$}&\resultplaceholder&\tablelang{ref.}{参照}&\resultplaceholder\\
&\tablelang{Scene flow / $G^\star$}{场景流/$G^\star$}&\resultplaceholder&\resultplaceholder&\resultplaceholder\\
&\tablelang{Sparse tracks / $G_0$}{稀疏轨迹/$G_0$}&\resultplaceholder&\resultplaceholder&\resultplaceholder\\
&\tablelang{Sparse tracks / $G^\star$}{稀疏轨迹/$G^\star$}&\resultplaceholder&\resultplaceholder&\resultplaceholder\\
\midrule
\tablelang{ID confirmation}{ID确认}&\tablelang{Selected complete design}{选定完整设计}&\resultplaceholder&\tablelang{ref.}{参照}&\resultplaceholder\\
&\tablelang{Matched action-only ablation}{匹配的仅动作消融}&\resultplaceholder&\resultplaceholder&\resultplaceholder\\
&\tablelang{Strongest complete alternative}{最强完整替代方案}&\resultplaceholder&\resultplaceholder&\resultplaceholder\\
\bottomrule
\end{tabularx}
\end{table*}


\paragraph{Rotation interface.}
The proprioception/action ablation compares rotation vectors with fixed-axis
XYZ Euler angles, using
$R=R_z(\gamma)R_y(\beta)R_x(\alpha)$,
$\beta\in[-\pi/2,\pi/2]$, and principal branches for the other angles.
Proprioception is absolute orientation; Action is a left-relative increment,
$R_{t+1}=\Delta R R_t$. Both encodings change together while physical frames,
commands, data, and budgets remain fixed. SR and SO(3) geodesic error provide
coordinate-independent evaluation. This is separate from the target-pose
readout's encoding.
